% Options for packages loaded elsewhere
% Options for packages loaded elsewhere
\PassOptionsToPackage{unicode}{hyperref}
\PassOptionsToPackage{hyphens}{url}
\PassOptionsToPackage{dvipsnames,svgnames,x11names}{xcolor}
%
\documentclass[
  letterpaper,
  DIV=11,
  numbers=noendperiod]{scrreprt}
\usepackage{xcolor}
\usepackage{amsmath,amssymb}
\setcounter{secnumdepth}{5}
\usepackage{iftex}
\ifPDFTeX
  \usepackage[T1]{fontenc}
  \usepackage[utf8]{inputenc}
  \usepackage{textcomp} % provide euro and other symbols
\else % if luatex or xetex
  \usepackage{unicode-math} % this also loads fontspec
  \defaultfontfeatures{Scale=MatchLowercase}
  \defaultfontfeatures[\rmfamily]{Ligatures=TeX,Scale=1}
\fi
\usepackage{lmodern}
\ifPDFTeX\else
  % xetex/luatex font selection
\fi
% Use upquote if available, for straight quotes in verbatim environments
\IfFileExists{upquote.sty}{\usepackage{upquote}}{}
\IfFileExists{microtype.sty}{% use microtype if available
  \usepackage[]{microtype}
  \UseMicrotypeSet[protrusion]{basicmath} % disable protrusion for tt fonts
}{}
\makeatletter
\@ifundefined{KOMAClassName}{% if non-KOMA class
  \IfFileExists{parskip.sty}{%
    \usepackage{parskip}
  }{% else
    \setlength{\parindent}{0pt}
    \setlength{\parskip}{6pt plus 2pt minus 1pt}}
}{% if KOMA class
  \KOMAoptions{parskip=half}}
\makeatother
% Make \paragraph and \subparagraph free-standing
\makeatletter
\ifx\paragraph\undefined\else
  \let\oldparagraph\paragraph
  \renewcommand{\paragraph}{
    \@ifstar
      \xxxParagraphStar
      \xxxParagraphNoStar
  }
  \newcommand{\xxxParagraphStar}[1]{\oldparagraph*{#1}\mbox{}}
  \newcommand{\xxxParagraphNoStar}[1]{\oldparagraph{#1}\mbox{}}
\fi
\ifx\subparagraph\undefined\else
  \let\oldsubparagraph\subparagraph
  \renewcommand{\subparagraph}{
    \@ifstar
      \xxxSubParagraphStar
      \xxxSubParagraphNoStar
  }
  \newcommand{\xxxSubParagraphStar}[1]{\oldsubparagraph*{#1}\mbox{}}
  \newcommand{\xxxSubParagraphNoStar}[1]{\oldsubparagraph{#1}\mbox{}}
\fi
\makeatother

\usepackage{color}
\usepackage{fancyvrb}
\newcommand{\VerbBar}{|}
\newcommand{\VERB}{\Verb[commandchars=\\\{\}]}
\DefineVerbatimEnvironment{Highlighting}{Verbatim}{commandchars=\\\{\}}
% Add ',fontsize=\small' for more characters per line
\usepackage{framed}
\definecolor{shadecolor}{RGB}{241,243,245}
\newenvironment{Shaded}{\begin{snugshade}}{\end{snugshade}}
\newcommand{\AlertTok}[1]{\textcolor[rgb]{0.68,0.00,0.00}{#1}}
\newcommand{\AnnotationTok}[1]{\textcolor[rgb]{0.37,0.37,0.37}{#1}}
\newcommand{\AttributeTok}[1]{\textcolor[rgb]{0.40,0.45,0.13}{#1}}
\newcommand{\BaseNTok}[1]{\textcolor[rgb]{0.68,0.00,0.00}{#1}}
\newcommand{\BuiltInTok}[1]{\textcolor[rgb]{0.00,0.23,0.31}{#1}}
\newcommand{\CharTok}[1]{\textcolor[rgb]{0.13,0.47,0.30}{#1}}
\newcommand{\CommentTok}[1]{\textcolor[rgb]{0.37,0.37,0.37}{#1}}
\newcommand{\CommentVarTok}[1]{\textcolor[rgb]{0.37,0.37,0.37}{\textit{#1}}}
\newcommand{\ConstantTok}[1]{\textcolor[rgb]{0.56,0.35,0.01}{#1}}
\newcommand{\ControlFlowTok}[1]{\textcolor[rgb]{0.00,0.23,0.31}{\textbf{#1}}}
\newcommand{\DataTypeTok}[1]{\textcolor[rgb]{0.68,0.00,0.00}{#1}}
\newcommand{\DecValTok}[1]{\textcolor[rgb]{0.68,0.00,0.00}{#1}}
\newcommand{\DocumentationTok}[1]{\textcolor[rgb]{0.37,0.37,0.37}{\textit{#1}}}
\newcommand{\ErrorTok}[1]{\textcolor[rgb]{0.68,0.00,0.00}{#1}}
\newcommand{\ExtensionTok}[1]{\textcolor[rgb]{0.00,0.23,0.31}{#1}}
\newcommand{\FloatTok}[1]{\textcolor[rgb]{0.68,0.00,0.00}{#1}}
\newcommand{\FunctionTok}[1]{\textcolor[rgb]{0.28,0.35,0.67}{#1}}
\newcommand{\ImportTok}[1]{\textcolor[rgb]{0.00,0.46,0.62}{#1}}
\newcommand{\InformationTok}[1]{\textcolor[rgb]{0.37,0.37,0.37}{#1}}
\newcommand{\KeywordTok}[1]{\textcolor[rgb]{0.00,0.23,0.31}{\textbf{#1}}}
\newcommand{\NormalTok}[1]{\textcolor[rgb]{0.00,0.23,0.31}{#1}}
\newcommand{\OperatorTok}[1]{\textcolor[rgb]{0.37,0.37,0.37}{#1}}
\newcommand{\OtherTok}[1]{\textcolor[rgb]{0.00,0.23,0.31}{#1}}
\newcommand{\PreprocessorTok}[1]{\textcolor[rgb]{0.68,0.00,0.00}{#1}}
\newcommand{\RegionMarkerTok}[1]{\textcolor[rgb]{0.00,0.23,0.31}{#1}}
\newcommand{\SpecialCharTok}[1]{\textcolor[rgb]{0.37,0.37,0.37}{#1}}
\newcommand{\SpecialStringTok}[1]{\textcolor[rgb]{0.13,0.47,0.30}{#1}}
\newcommand{\StringTok}[1]{\textcolor[rgb]{0.13,0.47,0.30}{#1}}
\newcommand{\VariableTok}[1]{\textcolor[rgb]{0.07,0.07,0.07}{#1}}
\newcommand{\VerbatimStringTok}[1]{\textcolor[rgb]{0.13,0.47,0.30}{#1}}
\newcommand{\WarningTok}[1]{\textcolor[rgb]{0.37,0.37,0.37}{\textit{#1}}}

\usepackage{longtable,booktabs,array}
\usepackage{calc} % for calculating minipage widths
% Correct order of tables after \paragraph or \subparagraph
\usepackage{etoolbox}
\makeatletter
\patchcmd\longtable{\par}{\if@noskipsec\mbox{}\fi\par}{}{}
\makeatother
% Allow footnotes in longtable head/foot
\IfFileExists{footnotehyper.sty}{\usepackage{footnotehyper}}{\usepackage{footnote}}
\makesavenoteenv{longtable}
\usepackage{graphicx}
\makeatletter
\newsavebox\pandoc@box
\newcommand*\pandocbounded[1]{% scales image to fit in text height/width
  \sbox\pandoc@box{#1}%
  \Gscale@div\@tempa{\textheight}{\dimexpr\ht\pandoc@box+\dp\pandoc@box\relax}%
  \Gscale@div\@tempb{\linewidth}{\wd\pandoc@box}%
  \ifdim\@tempb\p@<\@tempa\p@\let\@tempa\@tempb\fi% select the smaller of both
  \ifdim\@tempa\p@<\p@\scalebox{\@tempa}{\usebox\pandoc@box}%
  \else\usebox{\pandoc@box}%
  \fi%
}
% Set default figure placement to htbp
\def\fps@figure{htbp}
\makeatother


% definitions for citeproc citations
\NewDocumentCommand\citeproctext{}{}
\NewDocumentCommand\citeproc{mm}{%
  \begingroup\def\citeproctext{#2}\cite{#1}\endgroup}
\makeatletter
 % allow citations to break across lines
 \let\@cite@ofmt\@firstofone
 % avoid brackets around text for \cite:
 \def\@biblabel#1{}
 \def\@cite#1#2{{#1\if@tempswa , #2\fi}}
\makeatother
\newlength{\cslhangindent}
\setlength{\cslhangindent}{1.5em}
\newlength{\csllabelwidth}
\setlength{\csllabelwidth}{3em}
\newenvironment{CSLReferences}[2] % #1 hanging-indent, #2 entry-spacing
 {\begin{list}{}{%
  \setlength{\itemindent}{0pt}
  \setlength{\leftmargin}{0pt}
  \setlength{\parsep}{0pt}
  % turn on hanging indent if param 1 is 1
  \ifodd #1
   \setlength{\leftmargin}{\cslhangindent}
   \setlength{\itemindent}{-1\cslhangindent}
  \fi
  % set entry spacing
  \setlength{\itemsep}{#2\baselineskip}}}
 {\end{list}}
\usepackage{calc}
\newcommand{\CSLBlock}[1]{\hfill\break\parbox[t]{\linewidth}{\strut\ignorespaces#1\strut}}
\newcommand{\CSLLeftMargin}[1]{\parbox[t]{\csllabelwidth}{\strut#1\strut}}
\newcommand{\CSLRightInline}[1]{\parbox[t]{\linewidth - \csllabelwidth}{\strut#1\strut}}
\newcommand{\CSLIndent}[1]{\hspace{\cslhangindent}#1}



\setlength{\emergencystretch}{3em} % prevent overfull lines

\providecommand{\tightlist}{%
  \setlength{\itemsep}{0pt}\setlength{\parskip}{0pt}}



 


\usepackage{makeidx}
\makeindex
\KOMAoption{captions}{tableheading}
\makeatletter
\@ifpackageloaded{tcolorbox}{}{\usepackage[skins,breakable]{tcolorbox}}
\@ifpackageloaded{fontawesome5}{}{\usepackage{fontawesome5}}
\definecolor{quarto-callout-color}{HTML}{909090}
\definecolor{quarto-callout-note-color}{HTML}{0758E5}
\definecolor{quarto-callout-important-color}{HTML}{CC1914}
\definecolor{quarto-callout-warning-color}{HTML}{EB9113}
\definecolor{quarto-callout-tip-color}{HTML}{00A047}
\definecolor{quarto-callout-caution-color}{HTML}{FC5300}
\definecolor{quarto-callout-color-frame}{HTML}{acacac}
\definecolor{quarto-callout-note-color-frame}{HTML}{4582ec}
\definecolor{quarto-callout-important-color-frame}{HTML}{d9534f}
\definecolor{quarto-callout-warning-color-frame}{HTML}{f0ad4e}
\definecolor{quarto-callout-tip-color-frame}{HTML}{02b875}
\definecolor{quarto-callout-caution-color-frame}{HTML}{fd7e14}
\makeatother
\makeatletter
\@ifpackageloaded{bookmark}{}{\usepackage{bookmark}}
\makeatother
\makeatletter
\@ifpackageloaded{caption}{}{\usepackage{caption}}
\AtBeginDocument{%
\ifdefined\contentsname
  \renewcommand*\contentsname{Table of contents}
\else
  \newcommand\contentsname{Table of contents}
\fi
\ifdefined\listfigurename
  \renewcommand*\listfigurename{List of Figures}
\else
  \newcommand\listfigurename{List of Figures}
\fi
\ifdefined\listtablename
  \renewcommand*\listtablename{List of Tables}
\else
  \newcommand\listtablename{List of Tables}
\fi
\ifdefined\figurename
  \renewcommand*\figurename{Figure}
\else
  \newcommand\figurename{Figure}
\fi
\ifdefined\tablename
  \renewcommand*\tablename{Table}
\else
  \newcommand\tablename{Table}
\fi
}
\@ifpackageloaded{float}{}{\usepackage{float}}
\floatstyle{ruled}
\@ifundefined{c@chapter}{\newfloat{codelisting}{h}{lop}}{\newfloat{codelisting}{h}{lop}[chapter]}
\floatname{codelisting}{Listing}
\newcommand*\listoflistings{\listof{codelisting}{List of Listings}}
\makeatother
\makeatletter
\makeatother
\makeatletter
\@ifpackageloaded{caption}{}{\usepackage{caption}}
\@ifpackageloaded{subcaption}{}{\usepackage{subcaption}}
\makeatother
\makeatletter
\@ifpackageloaded{tcolorbox}{}{\usepackage[skins,breakable]{tcolorbox}}
\makeatother
\makeatletter
\@ifundefined{shadecolor}{\definecolor{shadecolor}{HTML}{f2f2f2}}{}
\makeatother
\makeatletter
\@ifundefined{codebgcolor}{\definecolor{codebgcolor}{HTML}{f2f2f2}}{}
\makeatother
\makeatletter
\ifdefined\Shaded\renewenvironment{Shaded}{\begin{tcolorbox}[sharp corners, enhanced, boxrule=0pt, breakable, borderline west={3pt}{0pt}{shadecolor}, frame hidden, colback={codebgcolor}]}{\end{tcolorbox}}\fi
\makeatother
\usepackage{bookmark}
\IfFileExists{xurl.sty}{\usepackage{xurl}}{} % add URL line breaks if available
\urlstyle{same}
\hypersetup{
  pdftitle={PyPedal --- Python Software for Pedigree Analysis},
  pdfauthor={John B. Cole},
  colorlinks=true,
  linkcolor={black},
  filecolor={Maroon},
  citecolor={Blue},
  urlcolor={blue},
  pdfcreator={LaTeX via pandoc}}


\title{PyPedal --- Python Software for Pedigree Analysis}
\author{John B. Cole}
\date{2025-12-04}
\begin{document}
\maketitle


\bookmarksetup{startatroot}

\chapter{Preface}\label{preface}

\section{Acknowledgments}\label{acknowledgments}

PyPedal was initially written to support the author's dissertation
research while at \href{http://www.lsu.edu/}{Louisiana State
University}. The initial development was supported in part by a grant
from The Seeing Eye, Inc. (Morristown, NJ). The author wishes to thank
Paul VanRaden for many helpful suggestions for improving the ability of
PyPedal to handle computations in large pedigrees. Additional feedback
in the form of bug reports, feature requests, and discussion of
computing strategies was provided by Dan Cieslak, Bradley J. Heins
(University of Minnesota), Edward H. Hagen (Institute for Theoretical
Biology, Humboldt-Universität zu Berlin), Kathy Hanford (University of
Nebraska), Thomas Kelly (Department of Animal and Poultry Science,
University of Guelph), Thomas von Hassell, Gianluca Saba, and Ken
Severin (University of Alaska -- Fairbanks). Gregor Gorjanc wrote a blog
entry describing
\href{http://ggorjan.blogspot.com/2007/04/installing-pypedal-under-debian.html}{how
to install PyPedal on Debian Linux}. Fernando Perez posted a LaTeX
preamble to the NumPy listserver that dramatically improved the PDF
version of this manual.

The Newfoundland pedigree was taken from the
\href{http://www.newfoundlanddog-database.net/en/}{NewFoundland Dog
database} and is used with permission.

The pedigree of European royalty used in the GEDCOM discussion,
\href{http://www.genealogyforum.com/gedcom/gedcom1/ged3.htm}{ged3.ged},
was taken from the \href{http://www.genealogyforum.com/}{Genealogy Forum
website}. It is believed to be in the public domain, and is used with
the knowledge of the website administrators.

\bookmarksetup{startatroot}

\chapter{Introduction}\label{introduction}

\begin{tcolorbox}[enhanced jigsaw, colframe=quarto-callout-note-color-frame, breakable, toprule=.15mm, leftrule=.75mm, rightrule=.15mm, colback=white, opacityback=0, arc=.35mm, left=2mm, bottomrule=.15mm]
\begin{minipage}[t]{5.5mm}
\textcolor{quarto-callout-note-color}{\faInfo}
\end{minipage}%
\begin{minipage}[t]{\textwidth - 5.5mm}

Any sufficiently disguised bug is indistinguishable from a feature. ---
Rich Kulawiec

\end{minipage}%
\end{tcolorbox}

\section{What is PyPedal?}\label{what-is-pypedal}

This chapter introduces PyPedal, a Python-language tool for working with
pedigrees.

PyPedal (\textbf{Py}thon \textbf{Ped}igree An\textbf{al}ysis) is a tool
for analyzing pedigree files. It calculates several quantitative
measures of genetic diversity from pedigrees, including average
coefficients of inbreeding and relationship, effective founder numbers,
and effective ancestor numbers. Checks are performed catch common
mistakes in pedigree files, such as parents with more recent birthdates
or smaller ID numbers than their offspring and animals appearing as both
sires and dams in the pedigree. Tools for pedigree visualization and
report generation are also provided. PyPedal only makes use of
information on pedigree structure, not individual genotypes. Genotypes
can be assigned to founders for use in gene-dropping simulations to
calculate the effective number of founder genomes.

PyPedal is a \href{http://www.python.org/}{Python} language module that
may be called by programs or used interactively. You must have Python
3.11 (or later) installed in order to use PyPedal. The NumPy module must
also be installed in order for you to use PyPedal. In addition, there
are several third-party packages used by PyPedal, which are discussed in
the chapter on installing PyPedal.

This manual is the official documentation for PyPedal. It includes a
tutorial and is the most authoritative source of information about
PyPedal with the exception of the source code. The tutorial material
will walk you through a series of manipulations of a simple pedigree.
All users of PyPedal are encouraged to follow the tutorial with a
working installation.

\section{Features}\label{features}

Key features of PyPedal include:

\begin{itemize}
\tightlist
\item
  Reading pedigree files in user-defined formats;
\item
  Checking pedigree integrity (duplicate IDs, parents younger than
  offspring, etc.);
\item
  Generating summary information such as frequency of appearance in the
  pedigree file;
\item
  Reordering and renumbering of pedigree files;
\item
  Computation of the numerator relationship matrix (\(A\)) from a
  pedigree file using the tabular method;
\item
  Inbreeding and relationship calculations for large pedigrees
  (pedigrees up to 600,000 animals have been processed with PyPedal);
\item
  Computation of coefficients of inbreeding using the Meuwissen and Luo
  (1992) and modified Meuwissen and Luo (Quaas, 1995, unpublished)
  algorithms;
\item
  Computation of average total and individual coefficients of inbreeding
  and relationship;
\item
  Calculation of coefficients of partial inbreeding using an iterative
  tabular method (Lacy, Alaks, and Walsh 1996; Gulisija et al. 2006);
\item
  Calculation of coefficients of ancestral inbreeding using the methods
  of Ballou (1997) or Suwanlee et al. (2007);
\item
  Decomposition of \(A\) into \(T\) and \(D\) such that \(A=TDT'\);
\item
  Computation of the direct inverse of \(A\) (not accounting for
  inbreeding) using the method of Henderson (1976);
\item
  Computation of the direct inverse of \(A\) (accounting for inbreeding)
  using the method of Quaas (1976);
\item
  Storage of \(A\) and its inverse between user sessions as persistent
  Python objects using the \texttt{pickle} module to avoid unnecessary
  calculations;
\item
  Calculation of theoretical and actual effective population sizes;
\item
  Computation of effective founder number using the exact algorithm of
  Lacy (1989);
\item
  Computation of effective founder number using the approximate
  algorithm of Boichard et al. (1997);
\item
  Computation of effective ancestor number using the algorithms of
  Boichard et al. (1997);
\item
  Selection of subpedigrees containing all ancestors of an animal;
\item
  Identification of the common relatives of two animals;
\item
  Calculation of the inbreeding of offspring from a prospective mating;
\item
  Output to ASCII text files, including matrices, coefficients of
  inbreeding and relationship, and summary information;
\item
  Importation and exportation of GEDCOM 5.5 files;
\item
  Importation and exportation of GENES 1.20 (dBase III) files;
\item
  Loading pedigrees from, and saving them to, MySQL, Postgres, and
  SQLite databases;
\item
  Simulation of pedigrees using an algorithm derived from that used in
  Matvec 1.1a;
\item
  Set operations on pedigrees (unions and intersections);
\end{itemize}

PyPedal has been used to perform calculations on pedigrees as large as
600,000 animals and has been used in scientific research (Cole, Franke,
and Leighton 2004; Cole 2007).

\section{Where to get information and
code}\label{where-to-get-information-and-code}

PyPedal and its documentation are available from the
\href{https://github.com/wintermind/pypedal3}{PyPedal 3 Github
repository}.

\bookmarksetup{startatroot}

\chapter{Getting Started}\label{getting-started}

For am example of how to load a pedigree, see
Listing~\ref{lst-load-pedigree}.

\begin{codelisting}

\caption{\label{lst-load-pedigree}}

\centering{

\begin{Shaded}
\begin{Highlighting}[]
\ImportTok{from}\NormalTok{ PyPedal }\ImportTok{import}\NormalTok{ pyp\_newclasses}
\NormalTok{options }\OperatorTok{=}\NormalTok{ \{}
    \StringTok{\textquotesingle{}renumber\textquotesingle{}}\NormalTok{: }\DecValTok{0}\NormalTok{,}
    \StringTok{\textquotesingle{}pedfile\textquotesingle{}}\NormalTok{: }\StringTok{\textquotesingle{}../test\_files/mrode.ped\textquotesingle{}}\NormalTok{,}
    \StringTok{\textquotesingle{}pedformat\textquotesingle{}}\NormalTok{: }\StringTok{\textquotesingle{}asd\textquotesingle{}}\NormalTok{,}
    \StringTok{\textquotesingle{}pedname\textquotesingle{}}\NormalTok{: }\StringTok{\textquotesingle{}Mrode Pedigree Table 2.1\textquotesingle{}}\NormalTok{,}
    \StringTok{\textquotesingle{}pedigree\_is\_renumbered\textquotesingle{}}\NormalTok{: }\DecValTok{1}\NormalTok{,}
    \StringTok{\textquotesingle{}messages\textquotesingle{}}\NormalTok{: }\StringTok{\textquotesingle{}quiet\textquotesingle{}}\NormalTok{,}
\NormalTok{\}}
\NormalTok{result }\OperatorTok{=}\NormalTok{ pyp\_newclasses.load\_pedigree(options)}
\NormalTok{result.metadata.printme()}
\end{Highlighting}
\end{Shaded}

\begin{verbatim}
Metadata for Mrode Pedigree Table 2.1 (../test_files/mrode.ped)
    Records:        6
    Unique Sires:       3
    Unique Dams:        2
    Unique Gens:        1
    Unique Years:       1
    Unique Founders:    2
    Unique Herds:       1
    Pedigree Code:      asd
\end{verbatim}

}

\end{codelisting}%

\bookmarksetup{startatroot}

\chapter*{References}\label{references}
\addcontentsline{toc}{chapter}{References}

\markboth{References}{References}

\phantomsection\label{refs}
\begin{CSLReferences}{1}{0}
\bibitem[\citeproctext]{ref-Ballou1997}
Ballou, J. D. 1997. {``Ancestral Inbreeding Only Minimally Affects
Inbreeding Depression in Mammalian Populations.''} \emph{Journal of
Heredity} 88: 169--78.

\bibitem[\citeproctext]{ref-ref352}
Boichard, D., L. Maignel, and E. Verrier. 1997. {``The Value of Using
Probabilities of Gene Origin to Measure Genetic Variability in a
Population.''} \emph{Genetics Selection Evolution} 29: 5--23.

\bibitem[\citeproctext]{ref-Cole2007a}
Cole, J. B. 2007. {``{P}y{P}edal: A Computer Program for Pedigree
Analysis.''} \emph{Computers and Electronics in Agriculture} 57:
107--13.

\bibitem[\citeproctext]{ref-Cole2004a}
Cole, J. B., D. E. Franke, and E. A. Leighton. 2004. {``Population
Structure of a Colony of Dog Guides.''} \emph{Journal of Animal Science}
82: 2906--12.

\bibitem[\citeproctext]{ref-GulisijaGWT2006}
Gulisija, D., D. Gianola, K. A. Weigel, and M. A. Toro. 2006.
{``Between-Founder Heterogeneity in Inbreeding Depression for Production
in {J}ersey Cows.''} \emph{Livestock Science} 104: 244--53.

\bibitem[\citeproctext]{ref-ref143}
Henderson, C. R. 1976. {``A Simple Method for Computing the Inverse of a
Numerator Relationship Matrix Used in Prediction of Breeding Values.''}
\emph{Biometrics} 32: 69--83.

\bibitem[\citeproctext]{ref-ref640}
Lacy, R. C. 1989. {``Analysis of Founder Representation in Pedigrees:
Founder Equivalents and Founder Genome Equivalents.''} \emph{Zoo
Biology.} 8: 111--23.

\bibitem[\citeproctext]{ref-LacyAW1996}
Lacy, R. C., G. Alaks, and A. Walsh. 1996. {``Hierarchical Analysis of
Inbreeding Depression in \emph{{P}eromyscus Polionotus}.''}
\emph{Evolution} 50: 2187--2200.

\bibitem[\citeproctext]{ref-ref585}
Meuwissen, H. E., and Z. Luo. 1992. {``Computing Inbreeding Coefficients
in Large Populations.''} \emph{Genetics Selection Evolution} 24:
305--13.

\bibitem[\citeproctext]{ref-ref235}
Quaas, R. L. 1976. {``Computing the Diagonal Elements and Inverse of a
Large Numerator Relationship Matrix.''} \emph{Biometrics} 32: 949--53.

\bibitem[\citeproctext]{ref-SuwanleeBSC2007}
Suwanlee, S., R. Baumung, J. Sölkner, and I. Curik. 2007. {``Evaluation
of Ancestral Inbreeding Coefficients: Ballou's Formula Versus Gene
Dropping.''} \emph{Conservation Genetics} 8: 489--95.

\end{CSLReferences}



\printindex


\end{document}
